\documentclass[11pt]{article}
\usepackage{amsmath,amssymb,amsthm}
\usepackage[margin=1.2in]{geometry}

\newtheorem{theorem}{Theorem}
\newtheorem{corollary}[theorem]{Corollary}
\newtheorem{definition}[theorem]{Definition}
\newtheorem{remark}[theorem]{Remark}
\newtheorem{conjecture}[theorem]{Conjecture}

\newcommand{\Deriv}{\mathsf{Deriv}}
\newcommand{\WC}{\mathsf{WC}}
\newcommand{\WCz}{\mathsf{WC}_0}
\newcommand{\SemTrue}{\mathsf{SemTrue}}
\newcommand{\EvalOK}{\mathsf{EvalOK}}
\newcommand{\ProofE}{\mathsf{ProofE}}
\newcommand{\evalCode}{\mathsf{evalCode}}
\newcommand{\thFunTFor}{\mathsf{thFunTFor}}
\newcommand{\substTFor}{\mathsf{substTFor}}
\newcommand{\reify}{\mathsf{reify}}
\newcommand{\code}{\mathsf{code}}
\newcommand{\Consistent}{\mathsf{Consistent}}
\newcommand{\Pair}{\mathsf{Pair}}
\newcommand{\TreeEq}{\mathsf{TreeEq}}
\newcommand{\BLevel}{\mathsf{BLevel}}
\newcommand{\enc}{\mathsf{enc}}
\newcommand{\NN}{\mathbb{N}}

\title{The Nelson--WC Boundary:\\
A Quantitative Analysis of Internal Soundness\\
in Recursive Arithmetic}
\author{}
\date{}

\begin{document}
\maketitle

\begin{abstract}
We formalize, in Agda with \texttt{--safe}, a quantitative boundary between
Nelson-style proof-by-proof internal soundness and G\"odel-style obstruction
in a system of recursive arithmetic over binary trees, following
Rose~\cite{rose1961}.
The system includes a 26-rule equational derivation calculus, a tree-recursive proof
checker ($\thFunTFor$), and a ground evaluator ($\evalCode$).
We define a well-chained certificate type $\WC$ with 16 constructors covering
all 26 rules, each carrying explicit local side conditions, and prove semantic
soundness: $\WC\;p \to \SemTrue\;p$.
Every genuine derivation yields a $\WC$ certificate once the relevant
$\EvalOK$ obligations are supplied.
We show that $\EvalOK$ cannot hold uniformly: an $\EvalOK$ instance for
the false equation $O = \Pair(O,O)$ implies inconsistency.
We identify a zero-backtracking fragment $\WCz$ where all side conditions
discharge automatically, and show this fragment includes the identity axiom,
reflexivity, symmetry, transitivity, unary congruence, and Schema~F.
The obstruction to universalization is localized to a specific
computational parameter---the backtracking level of semantic
verification---connecting Nelson's program in \emph{Elements}~\cite{nelson2015}
to the game-semantic complexity analysis of Aschieri~\cite{aschieri2016}.
\end{abstract}

%% ================================================================
\section{The System}

We work in a purely equational theory over binary trees
($\mathit{lf}$ and $\mathit{nd}\;a\;b$), with terms
$O$, $\mathit{var}\;n$, $\mathit{ap1}\;f\;t$, $\mathit{ap2}\;g\;t_1\;t_2$,
1-ary functors $I$, $\mathit{Fst}$, $\mathit{Snd}$, $\mathit{Comp}$,
$\mathit{Comp2}$, $\mathit{Rec}$, $\mathit{KT}$,
and 2-ary functors $\Pair$, $\mathit{Const}$, $\mathit{Lift}$,
$\mathit{Post}$, $\mathit{Fan}$, $\mathit{IfLf}$, $\TreeEq$.

The derivation system $\Deriv\;\mathit{hyp}\;\mathit{eq}$ is an
inductive type with 26 rules: 18 computation axioms (including
$\mathit{axI}$, $\mathit{axFst}$, $\mathit{axRecLf}$, $\mathit{axRecNd}$,
$\mathit{axTreeEqLL/LN/NL/NN}$, $\mathit{axIfLfL/N}$, etc.),
structural rules ($\mathit{axRefl}$, $\mathit{ruleSym}$,
$\mathit{ruleTrans}$, $\mathit{cong1}$, $\mathit{congL}$,
$\mathit{congR}$), variable instantiation ($\mathit{ruleInst}$),
hypothesis introduction ($\mathit{ruleHyp}$), and Schema~F
(uniqueness of tree recursion).

Schema~F says: if $f$ and $g$ both satisfy the same recursion
equation ($\mathit{Rec}\;z\;s$), then $f(x) = g(x)$ for all~$x$.
This replaces induction and is what makes the system strong
enough for G\"odel~II.

The system includes:
\begin{itemize}
\item A coding function $\code : \mathit{Term} \to \mathit{Tree}$
  with $\reify : \mathit{Tree} \to \mathit{Term}$.
\item A proof checker $\thFunTFor : \mathit{Fun1}$, defined as
  $\mathit{Rec}\;O\;\mathit{thFunStep}$, where $\mathit{thFunStep}$
  is a 26-branch dispatch.
\item A ground evaluator $\evalCode : \mathit{Fun1}$, defined as
  $\mathit{Rec}\;O\;\mathit{evalStep}$, dispatching on tag-$O$
  (returns $O$), tag-$\mathit{Ap1}$ (strips the functor), and
  default (passthrough).
\item A coded substitution $\substTFor : \mathit{Fun1}$, implementing
  variable replacement on coded terms.
\end{itemize}

Consistency is: $\Consistent\;\mathit{hyp} = \Deriv\;\mathit{hyp}\;(O = \Pair(O,O)) \to \bot$.

%% ================================================================
\section{G\"odel's Second Incompleteness Theorem}

Following Rose~\cite{rose1961}, we construct a G\"odel sentence via diagonal substitution.
The template equation is
\[
  \TreeEq(\thFunTFor(x),\;\substTFor(y)) = \Pair(O,O)
\]
with free variables $x$ (proof encoding) and $y$ (coded term).
Since $\TreeEq(a,b) = O$ iff $a = b$ and $\TreeEq(a,b) = \Pair(O,O)$ iff $a \ne b$,
this says: ``the equation proved by encoding~$x$ does not match the
coded substitution result.''

The G\"odel sentence $G$ is obtained by substituting the code of the
template itself for~$y$, yielding (after closing)
a fixed-point: $G$ says ``no proof encoding proves~$G$.''

\begin{theorem}[G\"odel II, \texttt{GodelII.agda}]
\[
  \ProofE\;\mathit{hyp} \;\to\; \Consistent\;\mathit{hyp} \;\to\;
  \Deriv\;\mathit{hyp}\;G \;\to\; \bot
\]
\end{theorem}

The proof uses $\mathit{thm14E}$ (every derivation has a correct proof encoding),
the diagonal construction (coded substitution closes to a fixed point),
and $\TreeEq(t,t) = O$ (proved via Schema~F, not metalevel computation).
From $\Deriv\;\mathit{hyp}\;G$ one derives both $\TreeEq(X,X) = \Pair(O,O)$
and $\TreeEq(X,X) = O$ for the same~$X$, giving $O = \Pair(O,O)$.

\begin{remark}
Schema~F was corrected during this development. The original formulation
required proving $f(\Pair(a,b)) = g(\Pair(a,b))$ without inductive hypothesis.
The corrected version requires $f$ and $g$ to satisfy the \emph{same} recursion
equation, matching Rose's original Schema~F (uniqueness of the fixed point of tree recursion).
\end{remark}

%% ================================================================
\section{The Well-Chained Fragment}

\begin{definition}
The \emph{semantic truth predicate} is:
\[
  \SemTrue(p, \mathit{hyp}) \;=\;
  \Deriv\;\mathit{hyp}\;
  (\evalCode(\mathsf{Fst}(\thFunTFor(p))) \;=\;
   \evalCode(\mathsf{Snd}(\thFunTFor(p))))
\]
The \emph{evaluation side condition} for $t = u$ is:
\[
  \EvalOK(t, u, \mathit{hyp}) \;=\;
  \Deriv\;\mathit{hyp}\;
  (\evalCode(\reify(\code(t))) \;=\; \evalCode(\reify(\code(u))))
\]
\end{definition}

\begin{definition}
$\WC\;p\;\mathit{hyp}$ is an inductive type with 16 constructors
covering all 26 derivation rules. Each constructor carries a
\emph{Nelson proof} (the $\thFunTFor$ reduction for that rule)
and \emph{local side conditions}.

The side conditions fall into three classes:
\begin{enumerate}
\item \textbf{None} ($\mathit{axRefl}$, $\mathit{ruleSym}$):
  semantic truth follows from the induction hypothesis or is trivial.
\item \textbf{Matching} ($\mathit{ruleTrans}$): the middle term of the
  two sub-proofs must agree.
\item \textbf{Evaluation} (base cases, $\mathit{ruleInst}$,
  $\mathit{congL/R}$): an $\EvalOK$-type condition that depends
  on the specific proof and cannot be derived from the IH.
\end{enumerate}

A generic constructor $\mathsf{wcBase}$ handles any rule given its
Nelson proof and evaluation condition.
\end{definition}

%% ================================================================
\section{The Nelson--WC Boundary}

\begin{theorem}[Local Soundness, \texttt{WellChained.agda}]
$\WC\;p\;\mathit{hyp} \;\Rightarrow\; \SemTrue(p, \mathit{hyp})$.
\end{theorem}

\begin{theorem}[Completeness with Side Conditions, \texttt{NelsonWCBoundary.agda}]
For every $d : \Deriv\;\mathit{hyp}\;(t = u)$ and hypothesis certificate
$\mathit{ph} : \ProofE\;\mathit{hyp}$, if $\EvalOK(t, u, \mathit{hyp})$
is supplied, then there exists $\enc(d)$ with $\WC\;\enc(d)\;\mathit{hyp}$.
\end{theorem}

\begin{theorem}[Obstruction, \texttt{NelsonWCBoundary.agda}]
\label{thm:obstruct}
$\EvalOK(O,\;\Pair(O,O),\;\mathit{hyp}) \;\Rightarrow\;
 \Consistent\;\mathit{hyp} \;\Rightarrow\; \bot$.
\end{theorem}

\begin{proof}
$\evalCode(\reify(\code(O))) = O$ by the tag-$O$ case.
$\evalCode(\reify(\code(\Pair(O,O))))$ is a $\Pair$: the code of $\Pair(O,O)$
has tag $\mathsf{tagAp2}$, which misses both tag-$O$ and tag-$\mathsf{Ap1}$ checks,
triggering passthrough.
So $\EvalOK$ gives $O = \Pair(A,B)$ for specific $A$, $B$.
Then $\TreeEq(O,O) = O$ and $\TreeEq(O, \Pair(A,B)) = \Pair(O,O)$
by the $\TreeEq$ axioms, giving $O = \Pair(O,O)$ by transitivity.
\end{proof}

\begin{corollary}[The Nelson--WC Boundary]
\label{cor:boundary}
\mbox{}
\begin{enumerate}
\item For each genuine derivation~$d$, the $\EvalOK$ obligation
  can be discharged (it reduces to a concrete computation on coded terms).
\item But $\EvalOK$ cannot hold uniformly, since it would then
  hold for false equations, contradicting consistency.
\item Therefore $\WC$ certificates and $\SemTrue$ cannot be
  internalized as provable universal statements.
\end{enumerate}
\end{corollary}

%% ================================================================
\section{The Zero-Backtracking Fragment}

\begin{definition}
$\WCz\;p\;\mathit{hyp}$ is a restricted inductive type with 7 constructors:
$\mathsf{wc0Refl}$, $\mathsf{wc0AxI}$, $\mathsf{wc0Sym}$,
$\mathsf{wc0Trans}$, $\mathsf{wc0TransO}$, $\mathsf{wc0Cong1}$,
$\mathsf{wc0RuleF}$.
No constructor carries external $\EvalOK$ obligations.
\end{definition}

The key lemma enabling this:

\begin{theorem}[Universal Functor Miss, \texttt{BoundaryQuantitative.agda}]
For every $\mathit{Fun1}$~$f$:
\begin{align*}
\TreeEq(\reify(\mathsf{codeF1}\;f),\; O) &= \Pair(O,O) \\
\TreeEq(\reify(\mathsf{codeF1}\;f),\; \mathsf{tagAp1T}) &= \Pair(O,O)
\end{align*}
Consequently, $\evalCode(\mathsf{mkAp1}(\mathsf{codeF1}\;f,\; t)) = \evalCode(t)$
for all~$f$ and all~$t$.
\end{theorem}

\begin{proof}
By case analysis on $f$ (7 cases: $I$, $\mathit{Fst}$, $\mathit{Snd}$,
$\mathit{Comp}$, $\mathit{Comp2}$, $\mathit{Rec}$, $\mathit{KT}$).
Each $\mathsf{codeF1}\;f$ has the form $\mathit{nd}(\mathsf{natCode}(26{+}k), \mathit{payload})$,
so $\reify$ gives $\Pair(\Pair(O, \ldots), \ldots)$, which misses both tag checks.
\end{proof}

\begin{theorem}[Zero-Backtracking Soundness, \texttt{BoundaryQuantitative.agda}]
\[
  \WCz\;p\;\mathit{hyp} \;\Rightarrow\; \SemTrue(p, \mathit{hyp})
\]
with \textbf{no external side conditions}.
\end{theorem}

\begin{proof}
By induction on $\WCz$. The $\mathsf{axI}$ and $\mathsf{axKT}$ cases use
the universal functor miss lemma. The $\mathsf{cong1}$ case uses
the fact that any $\Pair(\Pair(f_1,f_2),f_3)$ automatically passes
both miss checks (by $\mathsf{axTreeEqNL}$ and the $\mathsf{IfLf}$ argument).
The structural cases ($\mathsf{Sym}$, $\mathsf{Trans}$) compose from the
induction hypothesis.
\end{proof}

%% ================================================================
\section{The Quantitative Boundary}

The $\WCz$ fragment is the \emph{zero-backtracking} layer of a natural
stratification:

\begin{definition}
The \emph{backtracking level} $\BLevel$ of a $\WC$ certificate is:
\begin{itemize}
\item $\BLevel = 0$ for base cases and structural rules without $\mathit{ruleInst}$.
\item $\BLevel = k{+}1$ for $\mathsf{wcRuleInst}$ with a sub-proof of level~$k$.
\item $\BLevel = \max(\BLevel_1, \BLevel_2)$ for binary rules.
\end{itemize}
\end{definition}

The side condition for $\mathit{ruleInst}$ has the form:
\[
  \evalCode(\substTFor(L)) = \evalCode(\substTFor(R))
\]
where $\substTFor$ does not commute with $\evalCode$.
Each nested $\mathit{ruleInst}$ adds a layer of $\substTFor$
that must be verified independently.

The quantitative boundary:
\begin{center}
\begin{tabular}{lll}
$\BLevel = 0$ & free & automatic ($\WCz$, this paper) \\
$\BLevel = k$ & $k$-fold $\substTFor$ verification & per-instance \\
$\BLevel = \infty$ & impossible & Theorem~\ref{thm:obstruct}
\end{tabular}
\end{center}

Every concrete derivation has some finite $\BLevel$, but no single
internal theorem proves $\SemTrue$ for all levels simultaneously.

%% ================================================================
\section{Connection to Nelson's Program}

In \emph{Elements}~\cite{nelson2015}, Nelson proposes verifying
mathematical proofs ``from below'': for each specific proof, checking
that the needed bounds hold, rather than establishing them universally.
He emphasizes that the bounds depend on structural parameters (rank and
level) rather than on proof length.

Our formalization makes this precise:
\begin{itemize}
\item The $\WC$ type is the certificate format.
\item The explicit side conditions are Nelson's ``needed bounds.''
\item $\mathsf{semSound}$ is the local verification.
\item The obstruction theorem says exactly why the local strategy
  cannot be globalized.
\end{itemize}

Nelson's key insight in \emph{Elements} \S14 concerns the predicate
$\varepsilon(y) \leftrightarrow \forall x[x \uparrow y \text{ terminates}]$.
For each specific string~$X$, the system $S^*$ proves $\varepsilon(*X)$.
But $S^*$ cannot prove $\forall x\,\varepsilon(x)$.
The gap between ``for each specific~$X$'' and ``for all~$X$''
is exactly our gap between ``for each derivation, $\EvalOK$ is
satisfiable'' and ``$\EvalOK$ cannot hold uniformly.''
Nelson's $\varepsilon$ is the exponentiation-totality predicate;
our $\EvalOK$ is the evaluation-agreement predicate.
Both are local conditions that work proof-by-proof but resist universalization.

%% ================================================================
\section{Connection to Aschieri's Backtracking Analysis}

Aschieri~\cite{aschieri2016} shows that the complexity of
cut-elimination and strategy interaction is governed not by crude
proof depth but by the \emph{backtracking level}---the depth of
mind-change chains in the associated game-semantic strategy.
His Max-Min Theorem says: the real interaction complexity between
an $n$-backtracking and an $m$-backtracking strategy is controlled
by $\min(n,m)$, not $\max(n,m)$.

Our $\BLevel$ stratification has a strikingly parallel form:
\begin{itemize}
\item $\BLevel = 0$: internally manageable, no external obligations
  (like 1-backtracking strategies: safe in any context).
\item Complexity grows by controlled layers: each increment of
  $\BLevel$ corresponds to one extra $\substTFor$ verification.
\item The uniform ``all levels at once'' claim breaks
  (like trying to bound all backtracking levels simultaneously).
\end{itemize}

\begin{conjecture}[Nelson--Aschieri Connection]
There exists a finite-level hierarchy $\BLevel_k$ of semantic side
conditions such that:
\begin{enumerate}
\item $\BLevel_0$ corresponds to a zero-backtracking, obligation-free
  fragment ($\WCz$);
\item each increment of~$k$ corresponds to one extra layer of semantic
  backtracking through substitution;
\item every concrete derivation has some finite~$k$;
\item no single internal theorem proves $\BLevel_k$ for all~$k$.
\end{enumerate}
Moreover, $\BLevel_k$ is a Nelson-side semantic analogue of Aschieri's
game-semantic backtracking level.
\end{conjecture}

The salvaged version of Nelson's program is then:

\begin{quote}
Internal soundness certification is a stratified capability.
The system can verify itself at every finite backtracking level,
with costs controlled by proof structure, not proof length.
The inability to do it uniformly is the incompleteness theorem
in computational form.
\end{quote}

%% ================================================================
\section{What Is New}

The result is not a restatement of G\"odel~II.
G\"odel~II says: a consistent system cannot prove its own consistency.
We say something more specific:

\begin{quote}
The exact computational locus of incompleteness is the side condition
$\EvalOK$ needed to turn decoded proof outputs into semantic truth.
\end{quote}

The obstruction is not:
\begin{itemize}
\item diagonalization (the diagonal is constructed successfully);
\item self-reference ($\thFunTFor$ works correctly on all inputs);
\item recursion or complexity (all functions terminate; all proofs
  check in under 1 second);
\item proof encoding ($\ProofE$ exists for every derivation).
\end{itemize}

The obstruction \emph{is}: $\evalCode$, which evaluates coded terms
by tree-structural tag dispatch, cannot agree on the two sides of an
arbitrary coded equation. Concretely:
\begin{itemize}
\item $\evalCode$ dispatches on $\mathsf{Fst}$ of its input, conflating
  code structure with data structure ($\Pair(O,b)$ triggers tag-$O$).
\item $\substTFor$ does not commute with $\evalCode$: the
  $\mathit{ruleInst}$ side condition cannot be derived from the IH.
\item Any ground evaluator faces the same tension: to validate
  base-case axioms it must evaluate functors; but evaluating functors
  on coded terms requires distinguishing code-level tags from data-level
  structure, which cannot be done uniformly within the system's own
  equational theory.
\end{itemize}

%% ================================================================
\section{Formalization}

All results are in the \texttt{Guard/} directory, verified with
Agda~2.9.0 and \texttt{--safe} (no postulates, no \texttt{--type-in-type},
no \texttt{--with-K}).

\smallskip
\begin{center}
\begin{tabular}{ll}
\texttt{WellChained.agda} (475 lines) & $\WC$ type + $\mathsf{semSound}$ \\
\texttt{BoundaryQuantitative.agda} (293 lines) & $\WCz$ + $\mathsf{semSound0}$ \\
\texttt{NelsonWCBoundary.agda} (233 lines) & Bridge + obstruction \\
\texttt{Thm14E.agda} (1402 lines) & $\Deriv \to \ProofE$ \\
\texttt{GodelII.agda} & G\"odel~II \\
\texttt{NelsonDispatch.agda} (595 lines) & All 26 tag dispatches \\
\texttt{EvalCode*.agda} & Ground evaluator + lemmas \\
\texttt{Nelson*.agda} & Per-rule Nelson proofs (13+ tags) \\
\texttt{NelsonPipeline.agda} (277 lines) & Pipeline demonstrations \\
\end{tabular}
\end{center}

\begin{thebibliography}{9}
\bibitem{rose1961} H.\,E.\ Rose,
\emph{On the Consistency and Undecidability of Recursive Arithmetic},
Zeitschr.\ f.\ math.\ Logik und Grundlagen d.\ Math., Bd.~7, S.~124--135 (1961).

\bibitem{nelson1986} E.\ Nelson,
\emph{Predicative Arithmetic},
Mathematical Notes~32, Princeton University Press, 1986.

\bibitem{nelson2015} E.\ Nelson,
\emph{Elements},
arXiv:1510.00369 [math.LO], 2015 (unfinished, with afterword by S.\ Buss and T.\ Tao).

\bibitem{aschieri2016} F.\ Aschieri,
\emph{Game Semantics and the Geometry of Backtracking: a New Complexity
Analysis of Interaction},
arXiv:1511.06260v3 [math.LO], 2016.

\bibitem{kritchmanraz2010} S.\ Kritchman and R.\ Raz,
\emph{The Surprise Examination and the Second Incompleteness Theorem},
Notices of the AMS, \textbf{57}, 1454--1458, 2010.

\bibitem{guard1963} J.\ Guard,
\emph{Lecture Notes on Recursive Arithmetic}, 1963.
\end{thebibliography}

\end{document}
